% chapters/10/section-05.tex
\section{Program Pencarian Jalur Lengkap}
\label{sec:program-astar}

Bagian ini menyatukan seluruh unsur menjadi program pencarian jalur yang dapat
diamati dan diukur.

\subsection{Kode Lengkap}

\kodepython{python/ch10_astar.py}{Pencarian jalur A* pada peta berpetak beserta visualisasinya}{lst:ch10-astar}

\subsection{Membedah Susunan Programnya}

Program itu tersusun atas lima bagian yang jelas terpisah.

\begin{description}
  \item[Pembacaan peta] Fungsi \code{baca_peta} mengubah gambar huruf menjadi
        himpunan dinding, titik awal, dan titik tujuan. Fungsi itu juga
        memeriksa panjang baris yang harus sama.
  \item[Penentuan tetangga] Fungsi \code{tetangga} mengembalikan petak yang
        bersebelahan dan dapat dilalui, yaitu indikator \ind{a}.
  \item[Ukuran heuristik] Tiga fungsi kecil, yaitu \code{jarak_manhattan},
        \code{jarak_euclidean}, dan \code{tanpa_perkiraan}.
  \item[Pencarian] Fungsi \code{cari_jalur} menjalankan algoritma A* dengan
        antrean prioritas, lalu merekonstruksi jalurnya.
  \item[Pemeriksaan dan penggambaran] Fungsi \code{periksa_jalur} memeriksa
        keabsahan jalur, sedangkan fungsi \code{gambar_peta} menampilkan
        keadaan pencarian.
\end{description}

Pemisahan itu membuat setiap bagian dapat diganti tanpa mengganggu yang lain.
Mengubah bentuk peta, misalnya dari empat arah menjadi delapan arah, hanya
menuntut perubahan pada fungsi \code{tetangga}. Mengganti heuristik hanya
menuntut penambahan satu fungsi kecil.

\subsection{Keluaran Pengukuran}

Ketika program dijalankan dengan pilihan \code{--bandingkan}, terminal
menampilkan perbandingan ketiga heuristik pada peta yang sama.

\begin{kodekeluaran}
  Heuristik  Diperiksa   Jalur   Biaya  Temuan
--------------------------------------------------------------
  manhattan        274      51      50       0
  euclidean        353      51      50       0
      tanpa        437      51      50       0
--------------------------------------------------------------
Biaya jalur sama pada ketiga heuristik, sehingga ketiganya
menghasilkan jalur terpendek.
\end{kodekeluaran}

Kolom \istilah{Temuan} menyatakan banyaknya masalah yang ditemukan oleh
pemeriksaan jalur. Nilainya nol pada ketiga baris, yang berarti ketiga jalur
itu bersambung, dimulai dari titik awal, berakhir di titik tujuan, dan tidak
melewati dinding.

\subsection{Percobaan yang Disarankan}

\begin{table}[htbp]
\centering
\caption{Percobaan pada program pencarian jalur}
\label{tab:percobaan-ch10}
{\setlength{\tabcolsep}{4pt}\footnotesize%
\begin{tabularx}{\textwidth}{@{}L{5.6cm} Y@{}}
\toprule
\headrow \textbf{Yang diubah dan caranya} & \textbf{Yang sebaiknya Anda amati} \\
\midrule
Heuristik, melalui pilihan \code{--heuristic}, atau dengan menekan tombol H
ketika program berjalan. &
Bentuk penyebaran pencarian berubah: pada Manhattan penyebarannya sempit, pada
Euclidean lebih lebar, dan tanpa perkiraan menyebar ke segala arah. Jumlah
petak yang diperiksa bertambah, sedangkan biaya jalurnya tidak berubah. \\
\addlinespace
Titik tujuan, dengan memindahkan huruf \code{G} pada peta. &
Pencarian menyesuaikan arahnya. Apabila tujuannya terkurung dinding, jalur
kosong dan pemeriksaan melaporkan bahwa tujuan tidak dapat dicapai. \\
\addlinespace
Bentuk peta, dengan menambahkan dinding. &
Jalur berbelok mengelilinginya. Apabila dindingnya menutup seluruh jalan,
jalur menjadi kosong. \\
\addlinespace
Kecepatan penyajian, melalui pilihan \code{--langkah}. &
Nilai kecil membuat penyebaran pencarian lebih mudah diamati, sedangkan nilai
besar membuatnya selesai lebih cepat. \\
\addlinespace
Heuristik yang melebihkan jarak, yaitu dengan mengalikan hasil
\code{jarak_manhattan} dengan dua. &
Pencarian menjadi sangat cepat, tetapi biaya jalurnya dapat menjadi lebih besar
daripada 50. Inilah contoh heuristik yang tidak memenuhi syarat, dan
akibatnya terhadap optimalitas. \\
\bottomrule
\end{tabularx}}
\end{table}

\begin{peringatanbox}
Percobaan terakhir adalah yang paling berharga. Dengan mengalikan heuristik
dengan dua, pencarian berhenti lebih cepat dan memeriksa lebih sedikit petak,
tetapi jalur yang dihasilkan dapat lebih panjang. Program
tidak melaporkan kesalahan apa pun; satu-satunya tanda bahwa ada yang keliru
adalah angka biaya jalur yang lebih besar. Karena itu, biasakan memeriksa biaya
jalur, bukan hanya kecepatan pencarian.
\end{peringatanbox}
